\documentclass{article}

\usepackage{amsmath,amssymb}
\usepackage{graphicx}
\usepackage{xurl}
\usepackage[hidelinks]{hyperref}
\setlength{\emergencystretch}{2em}

% Shared commands for notes in docs/paper-gaps/.

\DeclareUnicodeCharacter{2080}{\ifmmode{}_{0}\else\textsubscript{0}\fi}
\DeclareUnicodeCharacter{2081}{\ifmmode{}_{1}\else\textsubscript{1}\fi}
\DeclareUnicodeCharacter{2082}{\ifmmode{}_{2}\else\textsubscript{2}\fi}
\DeclareUnicodeCharacter{2099}{\ifmmode{}_{n}\else\textsubscript{n}\fi}

\newcommand{\Lean}{\textsc{Lean}}
\ExplSyntaxOn
\NewDocumentCommand{\leanid}{m}{
  \ifmmode
    \textsf{\detokenize{#1}}
  \else
    \group_begin:
    \urlstyle{sf}
    \tl_set:Nn \l_tmpa_tl {#1}
    \tl_replace_all:Nnn \l_tmpa_tl {\_} {_}
    \exp_args:NV \path \l_tmpa_tl
    \group_end:
  \fi
}
\ExplSyntaxOff
\providecommand{\tr}{\operatorname{tr}}
\newcommand{\tnleanIssueBase}{https://github.com/LionSR/TNLean/issues/}
\newcommand{\tnleanPrBase}{https://github.com/LionSR/TNLean/pull/}
\newcommand{\tnleanDocBase}{https://sirui-lu.com/TNLean/docs/}
\newcommand{\ghissue}[1]{\href{\tnleanIssueBase#1}{GitHub issue \##1}}
\newcommand{\ghpr}[1]{\href{\tnleanPrBase#1}{PR \##1}}
\newcommand{\doclink}[1]{\href{\tnleanDocBase#1.html}{\path{#1}}}

% Dirac notation and the identity operator, as in blueprint/src/macros/common.tex.
% \providecommand keeps note-local definitions authoritative.
\usepackage{dsfont}
\providecommand{\ket}[1]{|#1\rangle}
\providecommand{\bra}[1]{\langle #1|}
\providecommand{\braket}[2]{\langle #1 | #2 \rangle}
\providecommand{\ketbra}[2]{|#1\rangle\!\langle #2|}
\providecommand{\Id}{\mathds{1}}
\providecommand{\one}{\mathds{1}}
\providecommand{\C}{\mathbb{C}}
\providecommand{\MN}[1]{M_{#1}(\C)}
\providecommand{\MD}{M_D(\C)}

% External referencing for paper sources.  The build helper writes sanitized
% auxiliary files under build/paper-aux/ and excludes duplicated source labels.
\usepackage{xr}

% Import a generated paper auxiliary file with a caller-supplied label prefix.
% The candidate paths support builds from docs/paper-gaps/, the repository root,
% and build/paper-aux/ itself.
\newcommand{\importpaperaux}[2]{%
  \IfFileExists{../../build/paper-aux/#2.aux}{%
    \externaldocument[#1]{../../build/paper-aux/#2}%
  }{%
    \IfFileExists{build/paper-aux/#2.aux}{%
      \externaldocument[#1]{build/paper-aux/#2}%
    }{%
      \IfFileExists{#2.aux}{%
        \externaldocument[#1]{#2}%
      }{}%
    }%
  }%
}

% General external reference macros
\newcommand{\paperref}[2]{\ref{#1-#2}}
\newcommand{\papereqref}[2]{(\ref{#1-#2})}

% CPSV16 source (1606.00608 / MPDO-22-12-17-2.tex) external document and shortcuts
\importpaperaux{cpsv16-}{1606.00608/MPDO-22-12-17-2-tnlean}
\newcommand{\cpsvref}[1]{\paperref{cpsv16}{#1}}
\newcommand{\cpsveqref}[1]{\papereqref{cpsv16}{#1}}

% Verdict marker: \gapnote{<kind>}{<status>}. Kind is the policy
% classification (clarification, local-correction, scope-restriction,
% unfaithful, false-source, open-gap); status is open, wip (elimination
% underway), resolved, or historical. Machine-read by the site index;
% prints a verdict line.
\newcommand{\gapnote}[2]{\par\noindent\textbf{Verdict:} #1 (#2).\par}

\importpaperaux{fbc25-}{2502.20257/main-tnlean}

\title{Source Boundary for MPU Action Tensors and Stabilizer Data}
\author{}
\date{2026-08-29}

\begin{document}
\maketitle
\gapnote{open-gap}{open}

\section{Scope and conclusion}

The target-paper coordinates in this note are pinned to
\path{Papers/2502.20257/main.tex}.  The three external source archives are not
duplicated in this repository.  They are pinned by exact version, stable
retrieval link, internal filename, and archive digest as follows.
\begin{itemize}
    \item The
          \href{https://arxiv.org/e-print/1706.07329v2}
          {arXiv:1706.07329v2 source archive} contains
          \path{cornerproblem.tex}; its SHA-256 digest is
          \nolinkurl{045a2a94589c6d4f82946bc9a1b4322ebd5109882355356191f462b3d362b23e}.
    \item The
          \href{https://arxiv.org/e-print/2203.12563v3}
          {arXiv:2203.12563v3 source archive} contains
          \path{REsubmission.tex}; its SHA-256 digest is
          \nolinkurl{53f0eb58d275d4bd801c9f38dc87e0928344745148860c97ba2e03fad4bd9b8c}.
    \item The
          \href{https://arxiv.org/e-print/2405.00439v2}
          {arXiv:2405.00439v2 source archive} contains
          \path{MPU-DW.tex}; its SHA-256 digest is
          \nolinkurl{166ce15a174167772f53d53b7ed15438fdf799717e053b128f612a264d3f8bb4}.
\end{itemize}
Every external line range below is therefore reproducible against the named
source file rather than a reformatted PDF.

The discussion of an MPU action on MPSs in
Ref.~\cite[lines~1782--1931]{FrancoRubio2025MPUGauging} uses two logically
different results.  The existence and scalar uniqueness of the action tensors
come from the MPS reduction theorem of
Ref.~\cite{Molnar2018SemiInjective}.  The description of scalar
$L$-symbols by stabilizer data comes from Sections~3 and~4.1 of
Ref.~\cite{GarreRubio2023MPOSym}.  The second paper also proves an existence
theorem for action tensors, but under an arbitrary-boundary invariance
hypothesis that is stronger than the periodic-boundary hypothesis in the
first paper and in Ref.~\cite{FrancoRubio2025MPUGauging}.

Thus there are two valid existence routes, but they cannot be interchanged.
For the simple-MPU setting of Ref.~\cite{FrancoRubio2025MPUGauging}, the
source-faithful route is the positive-length MPS reduction theorem, followed
by its nilpotency and scalar-uniqueness results.  Appendix~A of
Ref.~\cite{GarreRubio2023MPOSym} applies only after arbitrary-boundary
invariance has been supplied.  Once action and fusion tensors exist, the
scalar analysis of Section~4.1 of that paper applies independently of which
existence route produced them.

This distinction is the present open formal boundary.  The scalar
$L$-symbol relations are already formalized, whereas the tensor-valued
reduction and action data are not.

\section{The reduction route to action tensors}

Let $A_x$ be an injective MPS tensor for every $x$ in a set $\mathsf X$, and
let $B_{g,x}$ denote the MPS tensor obtained by stacking the MPU tensor
$\mathcal U_g$ on $A_x$ and vectorizing the physical legs.  The periodic
symmetry assumption is equality of the positive-length MPS vectors:
\begin{align}
    \mathcal V_N(B_{g,x})
        &=\mathcal V_N(A_{g\mathbin{\cdot}x}),
        &&N\geq 1.
    \label{eq:fbc25_action_audit_positive_length_symmetry}
\end{align}
Here the dot denotes the left action fixed below.  The target tensor
$A_{g\mathbin{\cdot}x}$ is injective and therefore normal.

Proposition~20 of Ref.~\cite[source lines~3124--3133, with proof at
lines~3815--3939]{Molnar2018SemiInjective} applies to a normal target tensor
$A$ and an arbitrary tensor $B$.  Its printed hypothesis allows an exponential
scalar $\lambda^N$, but its unscaled conclusion is false when $\lambda\ne1$.
The source proof establishes the equality specialization; the corrected
nonzero-$\lambda$ conclusion retains the factor $\lambda^n$ on words of
length $n$.  This source defect and its scalar counterexamples are recorded in
the
\href{https://sirui-lu.com/TNLean/paper-gaps/mgsc18_reduction_proportionality_scalar.pdf}
{dedicated paper-gap note}.  In the exact symmetry relation
(\ref{eq:fbc25_action_audit_positive_length_symmetry}), $\lambda=1$, so the
application here is unaffected.  The equality specialization supplies
matrices $V_{g,x}$ and $W_{g,x}$ such that
\begin{align}
    V_{g,x}W_{g,x}
        &=\mathbf 1,
    \label{eq:fbc25_action_audit_reduction_retraction}\\
    V_{g,x}B_{g,x}^{i_1}\cdots B_{g,x}^{i_n}W_{g,x}
        &=A_{g\mathbin{\cdot}x}^{i_1}\cdots
          A_{g\mathbin{\cdot}x}^{i_n},
        &&n\geq 1.
    \label{eq:fbc25_action_audit_reduction_words}
\end{align}
The proof obtains
(\ref{eq:fbc25_action_audit_reduction_retraction}) by taking the empty
interior word in the local reduction identity; it does not require equality
of the two length-zero periodic MPS coefficients.

The two matrices in
(\ref{eq:fbc25_action_audit_reduction_words}) are the source of the action
tensors $A^{\Gamma}_{g,x}$ and
$A^{\text{\reflectbox{$\Gamma$}}}_{g,x}$ in the target paper, after matching
the left and right virtual-leg orientations in its diagrams.  Proposition~21
and Definition~8 of Ref.~\cite[source lines~3142--3152, with proof at
lines~3943--3986]{Molnar2018SemiInjective} introduce the residual matrices
\begin{align}
    R_{g,x}^{i}
        &=B_{g,x}^{i}
          -W_{g,x}A_{g\mathbin{\cdot}x}^{i}V_{g,x}
    \label{eq:fbc25_action_audit_residual_definition}
\end{align}
and prove that their generated algebra is nilpotent.  The nilpotency length
$\ell_{g,x}$ is the least length after which every residual word vanishes:
\begin{align}
    R_{g,x}^{i_1}\cdots R_{g,x}^{i_n}
        &=0,
        &&n\geq \ell_{g,x}.
    \label{eq:fbc25_action_audit_residual_nilpotency}
\end{align}
Nilpotency alone does not remove the mixed terms in an expansion of $B$.
The additional input is the intended residual-word expansion of Appendix~B
and its preceding sandwiched-residual lemma
\cite[source lines~3946--3970 and 3993--4005]{Molnar2018SemiInjective}.
Suppress $g,x$ temporarily, write an empty product as the identity, and let
$\mathbf i=(i_1,\ldots,i_n)$.  The sandwiched residual words obey
\begin{align}
    V R^{i_1}\cdots R^{i_r}W
        &=0,
        &&r>0.
    \label{eq:fbc25_action_audit_sandwiched_residual_vanishes}
\end{align}

The first equality printed in Lemma~\texttt{B\_expand}, at source
lines~3993--4005, is false as printed and admits the following local
correction.  Its weak inequality
$0\leq r\leq s\leq n$ includes terms with an empty $A$-word and omits the
pure residual word.  At $n=1$ the printed right-hand side is
$WVR^i+WA^iV+R^iWV$, rather than $R^i+WA^iV$.  A nondegenerate reduction
counterexample and the correction are recorded in the dedicated false-source
note \cite{gap:mgsc18_residual_expansion_empty_word_error}.

Expanding $B^i=R^i+WA^iV$ and using
(\ref{eq:fbc25_action_audit_sandwiched_residual_vanishes}) shows that the
correct full expansion is
\begin{align}
    B^{i_1}\cdots B^{i_n}
        &=R^{i_1}\cdots R^{i_n}
          +\sum_{0\leq r<s\leq n}
          R^{i_1}\cdots R^{i_r}
          W A^{i_{r+1}}\cdots A^{i_s}V
          R^{i_{s+1}}\cdots R^{i_n}.
    \label{eq:fbc25_action_audit_full_residual_expansion}
\end{align}
Here the strict inequality $r<s$ says that the central $A$-word is nonempty.
Contracting this corrected identity and using residual nilpotency gives the
other two equalities printed in Lemma~\texttt{B\_expand}:
\begin{align}
    V B^{i_1}\cdots B^{i_n}
        &=\sum_{\max(0,n-\ell)\leq s\leq n}
          A^{i_1}\cdots A^{i_s}V
          R^{i_{s+1}}\cdots R^{i_n},
    \label{eq:fbc25_action_audit_left_residual_expansion}\\
    B^{i_1}\cdots B^{i_n}W
        &=\sum_{0\leq r\leq\min(\ell,n)}
          R^{i_1}\cdots R^{i_r}W
          A^{i_{r+1}}\cdots A^{i_n}.
    \label{eq:fbc25_action_audit_right_residual_expansion}
\end{align}
For $n<\ell$, the endpoints $s=0$ and $r=n$ in these contracted sums are the
pure residual terms $VR^{\mathbf i}$ and $R^{\mathbf i}W$.  For
$n\geq\ell$, the displayed endpoints with residual length exactly $\ell$
are zero by the source definition of nilpotency length.  They are retained to
match the printed bounds; the strictly surviving ranges are
$n-\ell<s\leq n$ and $0\leq r<\ell$.

Now let $\mathbf p$ and $\mathbf q$ both have length $m\geq\ell$.  The pure
residual term in the corrected full expansion of
$B^{\mathbf p\mathbf c\mathbf q}$ vanishes.  A central $A$-word which starts
after $\mathbf p$ leaves a residual prefix of length at least $\ell$, and one
which ends before $\mathbf q$ leaves a residual suffix of length at least
$\ell$.  Hence every surviving central $A$-word crosses all of $\mathbf c$.
The two contracted equations enumerate precisely these surviving terms and
give
\begin{align}
    B^{\mathbf p}W A^{\mathbf c}V B^{\mathbf q}
        &=B^{\mathbf p\mathbf c\mathbf q},
        &&|\mathbf p|=|\mathbf q|=m\geq\ell.
    \label{eq:fbc25_action_audit_exterior_from_residual_expansion}
\end{align}
This is the algebraic content of the exterior action identity
\papereqref{fbc25}{eq:action_exterior}; the interior identity
\papereqref{fbc25}{eq:action_interior} is the reduction relation
(\ref{eq:fbc25_action_audit_reduction_words}) and permits an interior length
$n\geq0$ \cite[lines~1787--1870]{FrancoRubio2025MPUGauging}.

Neither one-sided identity $VB^{\mathbf i}=A^{\mathbf i}V$ nor
$B^{\mathbf i}W=WA^{\mathbf i}$ follows.  For example,
$A^1=1$, $W=(1,0)^{\mathsf T}$, $V=(1,0)$, and
$B^1=\left(\begin{smallmatrix}1&1\\0&0\end{smallmatrix}\right)$ satisfy the
retraction and all-word reduction relations, while
$R^1=\left(\begin{smallmatrix}0&1\\0&0\end{smallmatrix}\right)$ has
$(R^1)^2=0$ but $V(B^1)^n=(1,1)\ne V$ for every $n>0$.  The full mixed-term
expansions, rather than residual nilpotency by itself, are therefore
load-bearing for
(\ref{eq:fbc25_action_audit_exterior_from_residual_expansion}).

Theorem~22 of Ref.~\cite[source lines~3154--3162, with proof at
lines~3990--4035]{Molnar2018SemiInjective} gives the precise uniqueness
statement.  If $(V,W)$ and $(\widetilde V,\widetilde W)$ are two reductions
whose nilpotency lengths are at most $N_0$, then there is
$c\in\mathbb C^\times$ such that, for every word of length $n>2N_0$,
\begin{align}
    VB^{i_1}\cdots B^{i_n}
        &=c\,\widetilde V B^{i_1}\cdots B^{i_n},
    \label{eq:fbc25_action_audit_reduction_left_uniqueness}\\
    B^{i_1}\cdots B^{i_n}W
        &=c^{-1}B^{i_1}\cdots B^{i_n}\widetilde W.
    \label{eq:fbc25_action_audit_reduction_right_uniqueness}
\end{align}
This is tailwise scalar uniqueness.  It is not the unconditional assertion
$V=c\widetilde V$ and $W=c^{-1}\widetilde W$ as bare matrices.

Theorem~1 of Ref.~\cite[source lines~358--377]{GarreRubioSchuch2024DomainWall}
restates the nilpotency-one case in the form used for MPU products and MPU
actions.  Its hypotheses give
\begin{align}
    VB^i\widehat V
        &=A^i,
    \label{eq:fbc25_action_audit_domain_wall_reduction}\\
    B^i\widehat V A^j V B^k
        &=B^iB^jB^k,
    \label{eq:fbc25_action_audit_domain_wall_exterior}\\
    V\widehat V
        &=\mathbf 1.
    \label{eq:fbc25_action_audit_domain_wall_retraction}
\end{align}
The statement is explicitly described there as informal: its footnote assumes
that the off-diagonal residual has nilpotency length one and says that this is
achieved by blocking.  Likewise, the target paper first permits
$m\geq\ell$ and then assumes that all nilpotency lengths have been reduced to
one by blocking
\cite[lines~1870--1874]{FrancoRubio2025MPUGauging}.  Consequently a formal
action-tensor theorem must either retain $\ell$ in its identities or prove the
blocking step before using the one-site identities.  Simplicity of the MPU is
not, by itself, a cited proof that an unblocked stacked tensor has
$\ell=1$.

The relevant global equality in the formal library is positive-length equality
for possibly different bond dimensions.  Full equality at length zero would
compare the traces of two identity matrices and therefore force equality of
the bond dimensions, although the stacked tensor $B_{g,x}$ and the reduced
tensor $A_{g\mathbin{\cdot}x}$ need not have equal bond dimensions.  Thus the
correct formal input is
\leanid{MPSTensor.SameMPV₂Pos}, not
\leanid{MPSTensor.SameMPV₂}.

\section{The arbitrary-boundary route is different}

Section~2 of Ref.~\cite[source lines~302--331]{GarreRubio2023MPOSym} defines
an MPS arbitrary-boundary subspace by allowing every virtual boundary matrix
$X$, and an MPO algebra by allowing every MPO boundary matrix $C$.  Its
closedness condition requires one boundary matrix, independent of the system
size, to represent each product of arbitrary-boundary MPOs.  Section~3 then
defines symmetry by the stronger condition that for every $C$ and $X$ there
is a boundary matrix $Y(C,X)$, again independent of $N$, such that
\begin{align}
    O^N_{T,C}\,\lvert\psi^N_{A,X}\rangle
        &=\lvert\psi^N_{A,Y(C,X)}\rangle,
        &&\text{for every }C,X,N.
    \label{eq:fbc25_action_audit_arbitrary_boundary_invariance}
\end{align}
This is Equation~(7) of that source
\cite[source lines~429--460]{GarreRubio2023MPOSym}.  Its Appendix~A starts
from the size-independent linear boundary map, takes a minimal-rank
decomposition, and uses the one-site and two-site identities plus block
injectivity to obtain fusion tensors and their orthogonality
\cite[source lines~2305--2454]{GarreRubio2023MPOSym}.  The final sentence
applies the same construction to
(\ref{eq:fbc25_action_audit_arbitrary_boundary_invariance}) to obtain action
tensors and their orthogonality
\cite[source line~2455]{GarreRubio2023MPOSym}.

By contrast, the simple-MPU hypothesis together with the independent periodic
symmetry assumption gives
(\ref{eq:fbc25_action_audit_positive_length_symmetry}) for the distinguished
closed boundaries.  Simplicity alone does not assert this equality.  These
hypotheses do not quantify over $C$ and $X$, and they do not supply the
boundary map used in Appendix~A.  The mismatch is structural, not merely a
missing lemma.  Ref.~\cite[source lines~1211--1216]{GarreRubio2023MPOSym}
proves that an injective group MPO satisfying $U_gU_h=U_{gh}$ and
$U_e=\mathbf 1$ must be on-site, with trivial $3$-cocycle, if its fusion
tensors also satisfy the arbitrary-boundary closedness condition.  Hence an
anomalous injective MPU representation with $U_e=\mathbf 1$ cannot satisfy
that condition.

There is a second difference.  In the arbitrary-boundary group algebra of
Ref.~\cite[source lines~638--683]{GarreRubio2023MPOSym}, the neutral operator
$O_e$ is generally a projector onto a relevant periodic-boundary subspace,
not the identity on the full Hilbert space, and the operators need not be
unitary.  The target paper instead uses simple injective MPUs, each unitary on
the full periodic chain, with $\mathcal U_e=\mathbf 1$.  The reduction theorem
of the preceding section is designed for this latter situation.

\section{Left-action and scalar conventions}

The target paper uses a left action: $g\mathbin{\cdot}x$ is the block obtained
by applying $g$ to $x$, and $gh\mathbin{\cdot}x$ means that $h$ acts first and
then $g$.  Its transitivity assumption is that for every $x,y\in\mathsf X$
there is $g\in G$ with $y=g\mathbin{\cdot}x$
\cite[lines~1782--1785]{FrancoRubio2025MPUGauging}.  With this convention,
the scalar compatibility equation is
\begin{align}
    L^x_{g,hk}L^x_{h,k}
        &=\omega(g,h,k)L^{k\mathbin{\cdot}x}_{g,h}L^x_{gh,k}.
    \label{eq:fbc25_action_audit_l_compatibility}
\end{align}
The shifted block is $k\mathbin{\cdot}x$, not
$g\mathbin{\cdot}x$ or $h\mathbin{\cdot}x$, because the comparison begins
with the rightmost action.  Equation
(\ref{eq:fbc25_action_audit_l_compatibility}) is
Equation~\papereqref{fbc25}{eq:omega_and_Ls} of the target paper
\cite[lines~1920--1924]{FrancoRubio2025MPUGauging} and agrees with the group
specialization in Ref.~\cite[source lines~683--719]{GarreRubio2023MPOSym}.

The target paper scales its reflected-$\Gamma$ action tensor by
$\gamma_{g,x}$ and its $\Gamma$ action tensor by $\gamma_{g,x}^{-1}$:
\begin{align}
    A^{\text{\reflectbox{$\Gamma$}}}_{g,x}
        &\longmapsto
          \gamma_{g,x}A^{\text{\reflectbox{$\Gamma$}}}_{g,x},
    &
    A^{\Gamma}_{g,x}
        &\longmapsto
          \gamma_{g,x}^{-1}A^{\Gamma}_{g,x}.
    \label{eq:fbc25_action_audit_action_tensor_gauge}
\end{align}
Together with its fusion-tensor convention this gives
\begin{align}
    L^x_{g,h}
        &\longmapsto
          \frac{\gamma_{gh,x}\,\beta_{g,h}}
               {\gamma_{g,h\mathbin{\cdot}x}\,\gamma_{h,x}}
          L^x_{g,h}.
    \label{eq:fbc25_action_audit_target_l_gauge}
\end{align}
These are Equations~\papereqref{fbc25}{eq:scalar_act_ten} and
\papereqref{fbc25}{eq:Lsymbgauge} of
Ref.~\cite[lines~1870--1919]{FrancoRubio2025MPUGauging}.
Equation~(19) of Ref.~\cite[source lines~720--726]{GarreRubio2023MPOSym}
prints the reciprocal factor.  The two formulas use inverse names for both
scalar gauges: replacing the 2023 parameters by
$\gamma^{-1}$ and $\beta^{-1}$ gives
(\ref{eq:fbc25_action_audit_target_l_gauge}).  The compatibility equation is
unchanged.  The formal library follows the target convention, so a proof may
not import the 2023 gauge formula without this inversion.

\section{Stabilizer data and the meaning of
  \texorpdfstring{$(H,\psi)$}{(H, psi)}}

Assume now that $G$ is finite, $\mathsf X$ is nonempty, and the left action is
transitive.  After choosing $x_0\in\mathsf X$, the subgroup in Section~4.1 of
Ref.~\cite[source lines~731--749]{GarreRubio2023MPOSym} is precisely the
stabilizer
\begin{align}
    H=\operatorname{Stab}_G(x_0)
        &=\{h\in G\mid h\mathbin{\cdot}x_0=x_0\}.
    \label{eq:fbc25_action_audit_stabilizer}
\end{align}
It is not an additional ``largest'' subgroup and it need not be normal.
Transitivity gives an equivariant bijection of sets
\begin{align}
    G/H&\longrightarrow\mathsf X,
    &gH&\longmapsto g\mathbin{\cdot}x_0,
    \label{eq:fbc25_action_audit_coset_equivalence}\\
    \lvert\mathsf X\rvert
        &=\lvert G/H\rvert.
    \label{eq:fbc25_action_audit_coset_cardinality}
\end{align}
Here $G/H$ is the left-coset set.  It is not a quotient group unless $H$ is
normal.  Replacing $x_0$ by $g\mathbin{\cdot}x_0$ replaces $H$ by the
conjugate subgroup $gHg^{-1}$, so the basepoint-free datum is a conjugacy
class of stabilizer data, not a distinguished subgroup.

Restricting
(\ref{eq:fbc25_action_audit_l_compatibility}) to $H$ shows that
$\omega|_H$ is cohomologically trivial.  In the convention printed in
Ref.~\cite[source lines~740--749]{GarreRubio2023MPOSym}, one chooses a
$2$-cochain $\beta$ satisfying
\begin{align}
    \omega(h_1,h_2,h_3)
    \frac{\beta_{h_1,h_2}\,\beta_{h_1h_2,h_3}}
         {\beta_{h_2,h_3}\,\beta_{h_1,h_2h_3}}
        &=1,
        &&h_1,h_2,h_3\in H,
    \label{eq:fbc25_action_audit_restricted_trivialization}
\end{align}
and then defines the representative
\begin{align}
    \psi_{x_0}(h_1,h_2)
        &=\frac{L^{x_0}_{h_1,h_2}}{\beta_{h_1,h_2}}.
    \label{eq:fbc25_action_audit_psi_representative}
\end{align}
In that chosen trivializing gauge, $\psi_{x_0}$ is a $2$-cocycle.  Changing
the action-tensor gauge changes it by a $2$-coboundary, so the invariant is its
class $[\psi_{x_0}]\in H^2(H,\mathbb C^\times)$.  Before choosing a
trivializer for $\omega|_H$, the corresponding object is a trivializing
$2$-cochain rather than a canonically based element of $H^2$.

Equation~(20) of Ref.~\cite{GarreRubio2023MPOSym} makes the
representative-level choices explicit; it is at lines~752--759 of the
official v3 source.  The prose at line~754 announces a solution for
$L^\alpha_{g_2,g_1}$, whereas the displayed equation at lines~755--758 is for
$L^\alpha_{g_1,g_2}$ and first resolves the $g_2$ action, then the $g_1$
action.  The prose has the swapped-order typo; the display fixes the convention
recorded here.  Choose a representative
$k_\alpha\in G$ for every coset and define $h_2,h_1\in H$ by
$g_2k_\alpha=k_\beta h_2$ and
$g_1k_\beta=k_\gamma h_1$.  The source gives
\begin{align}
    L^\alpha_{g_1,g_2}
        &=\frac{
          \omega(g_1g_2k_\alpha,h_2^{-1},h_1^{-1})\,
          \psi(h_1,h_2)}{
          \omega(g_1,g_2k_\alpha,h_2^{-1})\,
          \omega(g_1,g_2,k_\alpha)}.
    \label{eq:fbc25_action_audit_l_reconstruction}
\end{align}
Thus Equation~(20) constructs a chosen scalar $L$-symbol representative from
a chosen representative of $\omega$, a chosen trivialization and cocycle
representative, a basepoint, and coset representatives.  It is not a canonical
formula on cohomology classes, and it does not construct action tensors or an
MPU realization.

For fixed $(G,\omega,H)$, Ref.~\cite[source line~761]{GarreRubio2023MPOSym}
states that gauge classes of solutions form a torsor for
$H^2(H,\mathbb C^\times)$.  A choice of trivialization identifies that torsor
with the group, but in general there is no canonical solution corresponding
to the identity cohomology class.  Accordingly, the notation $(H,\psi)$ has
three levels which should not be conflated:
\begin{enumerate}
    \item in Equation~(20), $H$ and $\psi$ are chosen representatives, together
          with the unprinted auxiliary gauge choices just listed;
    \item for scalar classification, the invariant is stabilizer data up to
          basepoint conjugacy and a gauge class represented by $\psi$;
    \item for phases, the invariant is the resulting equivalence class of
          $L$-symbols, under the hypotheses of the phase theorem.
\end{enumerate}
The phase theorem itself concerns periodic Hamiltonians with exact MPS ground
spaces, equivalent $F$-symbols, and the restricted symmetric gapped paths
specified in that paper
\cite[source lines~1219--1272, especially the main result at line~1231]
{GarreRubio2023MPOSym}.  It is not a classification of arbitrary Hamiltonian
phases, raw cocycle representatives, action tensors, or MPU tensor data.
Conversely, the scalar construction
(\ref{eq:fbc25_action_audit_l_reconstruction}) proves existence of an
$L$-symbol solution for admissible chosen data, but not existence of tensor
data realizing every such solution in the simple-MPU setting.

\section{Present formal boundary}

The existing library separates the following layers.

\begin{itemize}
    \item TNLean defines
          \leanid{MPSTensor.IsPrimitivePaper}, the unnormalized uniform
          vector-spreading formulation of paper primitivity, and
          \leanid{MPSTensor.HasEventuallyFullKrausRank}.  The latter is
          equivalent to QICLean's \leanid{Kraus.IsNormal} by
          \leanid{MPSTensor.hasEventuallyFullKrausRank\_iff\_isNormal}.
          The reverse bridge from \leanid{MPSTensor.IsPrimitivePaper} to the
          word-span predicate in
          \leanid{MPSTensor.primitivePaper\_iff\_hasEventuallyFullKrausRank}
          assumes left-canonical normalization.  Thus an exact named formal
          predicate for source primitivity already exists, while any passage
          from it to word-span normality must retain its left-canonical
          normalization hypothesis.

          QICLean also defines \leanid{Kraus.IsInjective} and
          \leanid{Kraus.IsNBlkInjective} for one matrix family.  The latter
          says that the length-$N$ words of one tensor span its full matrix
          algebra.  None of these predicates is the paper's direct-sum
          block-injective condition
          \begin{align}
              A^i&=\bigoplus_{x\in\mathsf X}A_x^i,
              &\operatorname{span}\{A_x^i\}_i
                  &=\operatorname{Mat}_{D_x}(\mathbb C)
                  \quad\text{for every }x.
              \label{eq:fbc25_action_audit_block_injective_distinction}
          \end{align}
          The direct-sum assembly is available through
          \leanid{MPSTensor.toTensorFromBlocks} and the richer
          \leanid{MPSTensor.SectorDecomposition} infrastructure, but no
          existing predicate attaches a transitive group action and an MPU
          action to that family.

    \item TNLean defines \leanid{MPOTensor.IsMPU} by unitarity for every
          system size $N>1$, following the 2017 MPU convention, and defines
          \leanid{MPOTensor.IsMPUSimple} by the two local simple-tensor
          contractions.  It also has the product tensor
          \leanid{MPOTensor.mulTensor} and the doubled-index view
          \leanid{MPOTensor.toMPSTensor}.  It has no group-indexed MPU
          representation object and no tensor for the stacked action of an
          MPU on an MPS.  Those notions must precede the reduction theorem;
          they should not be inferred from the scalar $L$-symbol API.

    \item \leanid{MPSTensor.SameMPV₂Pos}
          expresses the global
          equality in
          (\ref{eq:fbc25_action_audit_positive_length_symmetry}) across
          different bond dimensions.  The missing result is the
          Moln\'ar--Ge--Schuch--Cirac reduction package:
          (\ref{eq:fbc25_action_audit_reduction_retraction})--
          (\ref{eq:fbc25_action_audit_residual_nilpotency}), its transport
          under blocking, and the tailwise uniqueness
          (\ref{eq:fbc25_action_audit_reduction_left_uniqueness})--
          (\ref{eq:fbc25_action_audit_reduction_right_uniqueness}).

    \item Mathlib's \leanid{MulAction},
          \leanid{MulAction.IsPretransitive},
          \leanid{MulAction.stabilizer}, and
          \leanid{MulAction.orbitEquivQuotientStabilizer} provide the correct
          action, transitivity, stabilizer, and coset-set language for
          (\ref{eq:fbc25_action_audit_stabilizer})--
          (\ref{eq:fbc25_action_audit_coset_equivalence}).  Transitivity is
          represented by pretransitivity together with nonemptiness.  No
          normality hypothesis on $H$ should be introduced.

    \item TNLean's \leanid{TNLean.Algebra.LSymbol} is the scalar function
          $\mathsf X\to G\to G\to\mathbb C^\times$.
          \leanid{TNLean.Algebra.LSymbol.IsCompatible} is exactly
          (\ref{eq:fbc25_action_audit_l_compatibility}), and
          \leanid{TNLean.Algebra.LSymbol.gauge} is exactly
          (\ref{eq:fbc25_action_audit_target_l_gauge}).  The name
          \leanid{TNLean.Algebra.ActionTensorGauge} denotes only the scalar
          cochain $\gamma_{g,x}$; it is not tensor data.  These declarations
          deliberately assert neither existence of action tensors nor their
          attachment to an MPU.

    \item \leanid{TNLean.Algebra.H2} is the quotient of genuine
          $\mathbb C^\times$-valued $2$-cocycles by coboundary equivalence.
          Its coefficient type is \leanid{Units}\,$\mathbb C$; despite prose in
          some module comments using $U(1)$, the implemented coefficient
          group matches the $\mathbb C^\times$ notation of the cited papers.
          No present theorem packages the restriction to a stabilizer, a
          chosen trivializer of $\omega|_H$, the torsor of solutions, or the
          representative reconstruction
          (\ref{eq:fbc25_action_audit_l_reconstruction}).

    \item The on-site, single-block specialization is already represented by
          \leanid{MPSTensor.IsOnSiteSymmetric} and
          \leanid{MPSTensor.virtual_rep_of_symmetric_injective}.  It should be
          reused when $H=G$ and the anomaly is trivial, but it is not a
          substitute for an anomalous MPU acting on several injective blocks.
          Likewise, the MPDO structure named
          \leanid{MPOTensor.BNTFusionTensorClause} concerns the CPSV16
          renormalization-fixed-point fusion clause and is not the group-MPU
          fusion tensor of the present sources.
\end{itemize}

The tensor-valued reduction package, rather than a new cohomology interface,
is therefore the next mathematical prerequisite.  Until it exists, the
formal scalar $L$-symbol results must not be described as proving the
existence or classification of MPU action tensors.

\section{Verdict}

The cited action-tensor claim is supported by
Ref.~\cite{Molnar2018SemiInjective}, provided that one retains the normal
target, positive-length periodic equality, residual nilpotency, and finite
blocking used by that theorem.  Its printed full residual expansion requires
the pure-residual and nonempty-middle correction in
(\ref{eq:fbc25_action_audit_full_residual_expansion}).  Its scalar uniqueness
is the tailwise
statement
(\ref{eq:fbc25_action_audit_reduction_left_uniqueness})--
(\ref{eq:fbc25_action_audit_reduction_right_uniqueness}).  The
arbitrary-boundary theorem of Ref.~\cite{GarreRubio2023MPOSym} is not an
alternative proof under the simple-MPU hypotheses.

The stabilizer classification is a classification of gauge classes of scalar
$L$-symbol solutions, and under the separate phase hypotheses it classifies
the corresponding exact-MPS symmetric phases.  The notation $(H,\psi)$ does
not by itself denote tensor data, and the solution set is naturally an
$H^2(H,\mathbb C^\times)$-torsor.  The formal gap remains open exactly at the
reduction-to-action-tensor layer; no new tensor data or cohomology interface is
introduced in this note.

\bibliographystyle{alpha}
\bibliography{references}

\end{document}
